\chapter{Fatigue, S-N curves, fracture mechanics}
\label{vol10:ch02}

\epigraph{Every cycle leaves a trace.}{}

\chapmeta{Half-life: durable. Archetypes: failure, scaling. Exercise target: 30. Project track: physical.}

\noindent
Fatigue is the failure of a material under repeated loading at stresses below its static strength. The phenomenon is responsible for a large fraction of mechanical-system failures: aircraft pressure cycles, bridge traffic cycles, rotating-machinery vibration, vehicle road loads, offshore platform wave loads. Fatigue failures often appear sudden but the damage accumulates over many cycles; the crack initiates, grows, and eventually causes brittle-like fracture even in ductile material. Fracture mechanics extends the analysis to components with pre-existing cracks (from manufacturing, in-service damage, or weld defects), predicting whether a crack will grow stably, propagate to instability, or arrest. This chapter develops fatigue's mechanism, S-N curves and their limits, variable-amplitude loading and Miner's rule, linear-elastic and elastic-plastic fracture mechanics, crack-growth analysis, inspection intervals, and the major historical fatigue cases.

\section{Fatigue mechanism: nucleation, propagation, fracture}\pathtag{core}

Fatigue failure proceeds in three stages:

\textbf{Nucleation:} cyclic stress produces microscopic slip in favourably-oriented grains; the slip bands intrude and extrude at the surface; persistent slip bands accumulate damage; a microcrack forms at the surface (or at a subsurface defect, inclusion, or stress concentration). The nucleation phase often consumes most of the fatigue life for smooth, ductile specimens; less so for components with stress raisers or defects.

\textbf{Propagation:} the microcrack grows under each subsequent load cycle. The growth rate is governed by the stress intensity at the crack tip; the crack typically extends a fraction of a micron per cycle in the early stages, accelerating as the crack grows. The crack surface shows characteristic markings (beach marks at the macroscale, striations at the microscale) that reveal the propagation history.

\textbf{Fracture:} when the crack reaches a critical size, the remaining cross-section can no longer support the load; the component fractures in a final overload. The fracture surface shows the smooth fatigue region (with beach marks) and a rougher overload region; the boundary defines the critical crack size.

The fatigue mechanism is independent of the static loading mode: tension, compression, bending, torsion, contact pressure, thermal cycling all produce fatigue. The driving stress is the cyclic component; the mean stress affects the rate but does not change the mechanism.

Fatigue is the dominant failure mode for many cyclically-loaded components: aircraft wings (gust and manoeuvre cycles), engine components (rotation and pressure cycles), rotating machinery (vibration), pressure vessels (pressurisation cycles), structures under wind or wave loading (millions of cycles per year). The discipline of fatigue-resistant design is essential.

\begin{archetype}[failure: as cumulative invisible damage]
Fatigue is the canonical cumulative-damage failure. Each cycle leaves a trace; the trace is invisible at first; the accumulated damage eventually crosses a threshold and the failure appears suddenly. The engineering response is recognition (the loading is cyclic; fatigue is a candidate failure mode), characterisation (the loading spectrum is known), design (the stress amplitude is below threshold for the design life), and inspection (the damage is detected before failure). The discipline applies broadly to other cumulative-damage modes: corrosion, wear, software-test-coverage degradation, accumulated operator stress.
\end{archetype}

\section{S-N curves and their limits}\pathtag{standard}

The S-N curve plots applied cyclic stress amplitude $S$ versus the number of cycles to failure $N$ for a given material, geometry, and loading. The curve is determined experimentally: specimens are loaded at various stress amplitudes; the number of cycles to failure is recorded; the data is fit to obtain the S-N curve.

The Wohler-Basquin equation: $\sigma_a = \sigma_f' (2N_f)^b$, where $\sigma_a$ is the stress amplitude, $N_f$ is the cycles to failure, $\sigma_f'$ is the fatigue strength coefficient (approximately the true fracture strength), and $b$ is the fatigue strength exponent (typically $-0.05$ to $-0.12$).

Steels and some other materials exhibit a fatigue limit (endurance limit): below a certain stress amplitude, the specimen does not fail in fatigue regardless of cycle count. The fatigue limit is typically $35\%$ to $50\%$ of the ultimate tensile strength for steels. Many non-ferrous metals (aluminum, copper, titanium) do not exhibit a fatigue limit; the curve continues to decline with cycle count, requiring design for a specified life rather than to a stress threshold.

The S-N curve is sensitive to many factors:

\textbf{Surface condition:} polished specimens have higher fatigue strength than as-machined or as-cast; surface roughness, scratches, and other surface defects reduce fatigue life. Surface treatments (shot peening, nitriding, carburising) introduce compressive residual stress and improve fatigue life.

\textbf{Stress concentration:} notches, holes, fillets, and other geometry concentrate stress at the location of nucleation. The fatigue strength is reduced by the stress concentration factor (often nominally; the actual reduction is given by the fatigue notch factor $K_f$, which is somewhat less than the elastic concentration factor $K_t$).

\textbf{Mean stress:} tensile mean stress reduces fatigue life; compressive mean stress can extend it. The Goodman, Gerber, and Soderberg relationships model the effect.

\textbf{Environment:} corrosive environments dramatically reduce fatigue life (corrosion fatigue). The fatigue limit may disappear in some environments. Saltwater, humid air, and chemical processes are common environmental fatigue concerns.

\textbf{Temperature:} elevated temperatures reduce fatigue strength; creep-fatigue interaction at high temperature is a separate failure mode for turbine and reactor components.

\textbf{Specimen size:} larger specimens have lower fatigue strength because of the larger volume of material with possible defects; the size effect is logarithmic with volume.

The S-N curve in practice: design specifications use the curve appropriate to the actual conditions (with margins for the uncertainties). Standards (ASTM E466, ASTM E467) define the testing procedures.

\section{Variable amplitude, Miner's rule}\pathtag{standard}

Real components rarely see constant-amplitude loading; the load varies over time. Miner's linear damage rule (Palmgren-Miner) is the standard simple model: each load cycle at amplitude $S_i$ consumes a fraction $n_i/N_i$ of the fatigue life, where $N_i$ is the cycles to failure at $S_i$ from the S-N curve; the component fails when $\sum n_i/N_i = 1$.

Miner's rule is approximate. The actual damage depends on the order of loading (sequence effects), the interaction between high and low amplitudes, and the presence of overloads (which can retard subsequent crack growth via compressive residual stress at the crack tip). The rule's accuracy is generally within a factor of $2$ to $3$ for many engineering applications; the simplicity makes it widely used.

Load-spectrum analysis is the discipline of characterising the cyclic loading. The standard procedures:

\textbf{Rainflow counting:} decomposes a complex load history into individual cycles, each with its own amplitude and mean stress. The rainflow algorithm (developed in the 1960s and now standard in fatigue software) handles arbitrary sequences and is the basis for variable-amplitude analysis.

\textbf{Load spectrum:} the histogram of cycle amplitudes for the expected service life. Aircraft load spectra are documented by aircraft type and mission profile; vehicle load spectra by road type and use; offshore-structure spectra by wave-climate models.

\textbf{Damage calculation:} the rainflow-decomposed cycles are applied via Miner's rule (or a more sophisticated model) to compute the predicted life.

\textbf{Block-loading testing:} the predicted spectrum is applied to physical specimens to verify the calculation. Service-life simulation rigs replicate the loading on representative specimens.

Aircraft fatigue analysis (the FALSTAFF and equivalent standardised load sequences) and ground-vehicle fatigue analysis are mature disciplines. The methods inform the design (fatigue life of the component), the inspection intervals, and the service-life limits.

\section{Linear-elastic fracture mechanics}\pathtag{standard}

Linear-elastic fracture mechanics (LEFM) analyses the behaviour of components with cracks under elastic loading. The stress field near a crack tip is singular; the magnitude of the singularity is described by the stress intensity factor $K$.

For a crack of length $2a$ in an infinite plate under remote tensile stress $\sigma$, the stress intensity is $K_I = \sigma \sqrt{\pi a}$. The subscript $I$ denotes Mode I (opening); Modes II (sliding) and III (tearing) describe other crack-opening modes; superposition gives the combined $K$.

The fracture toughness $K_{Ic}$ is the material property: the critical $K$ above which the crack extends unstably. $K_{Ic}$ is measured by standardised tests (ASTM E399) on compact-tension or single-edge-notch-bend specimens. Typical values: steel $50$ to $200\,\text{MPa}\sqrt{\text{m}}$; aluminum $20$ to $50$; titanium $50$ to $100$; ceramics $1$ to $10$; concrete $0.3$ to $0.5$.

The fracture criterion: failure occurs when $K_I \geq K_{Ic}$. For a given material ($K_{Ic}$) and component geometry (function relating $K_I$ to crack size and applied load), the critical crack size for fracture is calculable. The smaller of $K_{Ic}$ (linear-elastic) and $K_c$ (plain-strain vs plane-stress correction) is used.

LEFM applies when the plastic zone at the crack tip is small compared to the crack length and the component dimensions. This applies to brittle materials, high-strength materials, and thick sections. For ductile materials in plane-stress conditions or with large plastic zones, elastic-plastic fracture mechanics (next section) is more appropriate.

The J-integral provides an energy-based criterion: the path-independent J-integral around the crack tip equals the energy release rate; failure occurs when $J \geq J_{Ic}$. The criterion handles some nonlinear cases that the K-based criterion does not.

\section{Elastic-plastic fracture mechanics}\pathtag{mastery}

Elastic-plastic fracture mechanics (EPFM) handles cases where the plastic zone at the crack tip is significant. The standard parameters:

\textbf{Crack-tip opening displacement (CTOD):} the opening of the crack at its tip due to plastic deformation. Critical CTOD ($\delta_c$) is the fracture criterion. Used in plant-construction codes (BS 7910, API 579).

\textbf{J-integral:} extended to elastic-plastic behaviour with material-specific J-resistance curves. The J-R curve describes how J increases with crack extension; the tearing modulus quantifies the slope.

\textbf{Tearing instability:} crack extends stably under increasing load until the tearing modulus exceeds a critical value; then extends unstably. The instability criterion is critical for pressure vessels, where leak-before-break is the design goal (the crack extends stably until it penetrates the wall, leaking at low pressure rather than rupturing catastrophically).

EPFM is computationally more complex than LEFM (the J-integral requires numerical computation in many practical geometries) but is required for ductile materials, thin sections, and components operating in the plastic regime.

The leak-before-break philosophy is the engineering goal for many pressure vessels: design the wall thick enough and the material tough enough that any crack would penetrate the wall (causing detectable leakage) before reaching the size at which catastrophic rupture would occur. The discipline is mature in nuclear-pressure-vessel design.

\section{Crack-growth analysis}\pathtag{standard}

Crack-growth analysis predicts how a crack of given initial size will grow under cyclic loading. The Paris law: $\frac{da}{dN} = C (\Delta K)^m$, where $\Delta K = K_{max} - K_{min}$ is the stress intensity range and $C$, $m$ are material constants. Typical $m$ values: $2.5$ to $4$ for metals; the constants depend on environment.

Integrating the Paris law from initial crack size $a_0$ to final size $a_f$ gives the number of cycles to failure. The integration accounts for the changing $\Delta K$ as the crack grows; numerical integration is standard.

Real cracks deviate from the Paris law in the threshold regime (below $\Delta K_{th}$, growth is negligible) and the upper regime (approaching $K_{Ic}$, growth accelerates rapidly). The Forman, NASGRO, and other equations capture these regimes with additional constants.

Load-interaction effects modify the crack-growth rate: overloads can retard subsequent growth (compressive residual stress at the crack tip); underloads can accelerate growth (relaxation of the retardation effect). The Willenborg, Wheeler, and contact-stress models account for these effects.

Aircraft damage-tolerance design uses crack-growth analysis to set inspection intervals: assume a flaw of detectable size; compute the cycles to grow to critical size; inspect at intervals smaller than this time. The discipline is mature in commercial aviation (FAR 25.571) and in military aviation (MIL-A-83444).

\section{Inspection intervals and damage tolerance}\pathtag{standard}

Damage-tolerance design accepts that flaws exist and ensures that they will be detected before they cause failure. The design and operating discipline includes:

\textbf{Initial flaw assumption:} assume a flaw of size $a_i$ exists at delivery (from manufacturing). The size depends on the inspection capability at acceptance.

\textbf{Growth analysis:} compute the cycles for the flaw to grow from $a_i$ to the critical size $a_c$. The analysis uses the load spectrum, the crack-growth law, and the geometry's stress-intensity expression.

\textbf{Inspection interval:} set the inspection interval to a fraction of the growth life (typically $1/3$ to $1/2$). At each inspection, the inspector verifies that no flaw has grown above the detectability threshold.

\textbf{Detection capability:} the inspection method has a probability-of-detection (POD) curve as a function of flaw size. Below the threshold, detection is unlikely; above, reliable. The detection threshold sets the size $a_i$ assumed for the next inspection interval.

\textbf{Action on detection:} if a flaw is detected, the component is repaired (e.g., by replacement or splice), retired, or operated at reduced load with shorter inspection intervals.

Commercial aviation's damage-tolerance design is the canonical example: the Boeing 737 fuselage, the Airbus A320 wing, etc., are designed with assumed initial flaws, calculated growth, and scheduled inspection. The Aloha Airlines Flight 243 accident (\cite{acc:ntsb-aar94-04}, 1988) was the key event that drove the industry from safe-life to damage-tolerance philosophy; the multi-site fatigue cracking in the aged fuselage had not been adequately inspected.

\section{Case: de Havilland Comet, Aloha 243, Liberty Ships}\pathtag{core}

Three historical fatigue cases illustrate the discipline's importance:

\textbf{de Havilland Comet (1954):} the first commercial jet airliner, entered service in 1952. Two Comets suffered in-flight breakups in 1954 (BOAC Flight 781 over Elba in January; SAA Flight 201 over Naples in April). The investigation (\cite{acc:cohen-comet-report-1955}, \cite{text:withey-fatigue-comet}) identified fatigue cracking at the corners of the rectangular windows; the stress concentration at the corners, combined with the pressurisation cycles, led to crack initiation and growth. The cracks propagated to a critical size; the fuselage rupture was catastrophic.

The Comet investigation established several principles that became standard:

\textbf{Pressure cycling fatigue:} aircraft pressurisation cycles are a significant fatigue load. The Comet's fuselage had been designed for static pressure, not for the cyclic fatigue.

\textbf{Stress-concentration design:} sharp corners are stress concentrators; the rectangular windows were replaced by oval windows; subsequent aircraft have rounded windows for this reason.

\textbf{Full-scale fatigue testing:} aircraft are tested by simulating the entire service load cycle on a complete structure. The Comet's redesigned fuselage was tested in this way; the practice is now standard.

\textbf{Aloha Airlines Flight 243 (1988):} a Boeing 737-200 lost a $5\,\text{m}$ section of the upper fuselage at $24{,}000\,\text{ft}$; one flight attendant was lost; the aircraft landed safely with extensive damage. The investigation (\cite{acc:ntsb-aar94-04}) identified multi-site fatigue cracking at the rivet rows; the cracks had initiated at multiple locations along the same rivet row; they grew and joined into a continuous crack that ran the length of the section. The aircraft was $20$ years old with $89{,}680$ cycles, well above the original design life.

The Aloha investigation drove:

\textbf{Aging-aircraft initiative:} the FAA and industry recognised that aircraft beyond original design life require specific maintenance and inspection programs.

\textbf{Damage-tolerance philosophy:} the discipline of damage-tolerance design (with explicit acknowledgement of crack growth) became mandatory.

\textbf{Multi-site damage:} the design must consider that multiple cracks may initiate at similar locations and interact.

\textbf{Inspection enhancement:} non-destructive inspection methods improved to detect smaller cracks reliably.

\textbf{Liberty Ships (1940s):} during World War II, the US built $2{,}710$ Liberty cargo ships using welded construction (a new technology for ships at the scale). Approximately $200$ Liberty Ships suffered serious hull cracking; $19$ broke in half. The investigation identified brittle fracture: the welded hull (unlike riveted construction, where cracks would stop at the rivet holes) propagated cracks across the entire structure; the steel had inadequate toughness at the cold ocean temperatures; stress concentrations at hatch corners initiated the cracks.

The Liberty Ship investigation drove:

\textbf{Fracture mechanics:} the discipline emerged in part from the need to understand the Liberty Ship failures; the field formalised in the 1950s.

\textbf{Steel-toughness specification:} the Charpy V-notch impact test became standard for specifying toughness at operating temperature; the ductile-to-brittle transition temperature is a design parameter.

\textbf{Crack arrest:} the design of crack-stopping features (rivet rows, beams, crack-arresting strakes) became standard for large welded structures.

These three cases established the modern fatigue and fracture discipline; the lessons remain in every aircraft, pressure vessel, ship, and large structure built today.

\begin{principle}[The damage-tolerance principle]
For components subject to cyclic loading, the design must acknowledge that cracks will initiate and grow. The discipline is: characterise the loading spectrum, compute the crack-growth life, set inspection intervals smaller than the growth life, qualify inspections for the required detection threshold, plan the repair or replacement when cracks are found. The safe-life philosophy (no flaws allowed) is preferred in some specific applications but is unreliable in general; the damage-tolerance philosophy (flaws expected and managed) is the engineering default for most cyclic-loaded structures.
\end{principle}

\section{Project}\pathtag{core}

\begin{project}[Re-analyse the Aloha 243 fatigue case from primary FAA documents]
\textbf{Objective.} Re-analyse the Aloha Airlines Flight 243 fatigue case using the primary FAA and NTSB investigation documents. Identify the decision points where different choices would have prevented the failure. Develop the engineer's own assessment of the lessons.

\textbf{Method.} Read the NTSB accident report (\cite{acc:ntsb-aar94-04}), the FAA's airworthiness directives following the accident, and (if available) the post-accident industry studies on aging aircraft. Identify the engineering decisions: original design assumptions about pressure-cycle fatigue, inspection programs and intervals, NDI methods, maintenance practice, retirement criteria. For each decision, identify whether a different choice would have detected or prevented the failure; quantify the cost-benefit trade-offs that were involved.

\textbf{Deliverables.} The primary sources reviewed; the timeline of design and operational decisions; the engineering analysis of each decision point; your synthesis of the lessons that generalise beyond this specific aircraft.

\textbf{Hazard.} The investigation involves loss of life; engage with the human reality alongside the engineering analysis.
\end{project}

\section{Exercises}

\subsection*{Calculation}\pathtag{core}

\begin{exercise}
A steel specimen has fatigue strength $\sigma_f' = 1200\,\text{MPa}$ and fatigue exponent $b = -0.085$. Compute the cycles to failure at stress amplitude $400\,\text{MPa}$.
\end{exercise}

\begin{exercise}
A component is loaded with cycles at three stress amplitudes: $200\,\text{MPa}$ ($10^5$ cycles, $N_f = 10^7$), $300\,\text{MPa}$ ($10^4$ cycles, $N_f = 5 \times 10^5$), $400\,\text{MPa}$ ($10^3$ cycles, $N_f = 10^4$). Apply Miner's rule to determine fraction of life used.
\end{exercise}

\begin{exercise}
A through-thickness crack of length $2a = 10\,\text{mm}$ exists in an infinite plate of steel ($K_{Ic} = 80\,\text{MPa}\sqrt{\text{m}}$). Compute the critical applied stress for fracture.
\end{exercise}

\begin{exercise}
The Paris law constants for an aluminum alloy are $C = 10^{-11}$ (units $\text{m/cycle}$, $\text{MPa}\sqrt{\text{m}}$), $m = 3.0$. Compute the cycles to grow a crack from $0.5\,\text{mm}$ to $5\,\text{mm}$ at constant $\Delta K = 10\,\text{MPa}\sqrt{\text{m}}$.
\end{exercise}

\begin{exercise}
An aircraft fuselage has pressurisation cycles of $\Delta \sigma = 100\,\text{MPa}$ at frame stations. A crack of initial size $a_0 = 1\,\text{mm}$ grows by Paris law with $C = 10^{-12}$, $m = 3.5$. Compute the cycles to reach $a_c = 25\,\text{mm}$, the critical size for the panel.
\end{exercise}

\begin{exercise}
A steel with $\sigma_{UTS} = 800\,\text{MPa}$ has fatigue limit at $0.5 \sigma_{UTS}$. The stress amplitude at a notched location is $300\,\text{MPa}$ with stress concentration $K_t = 2.5$. State whether infinite life is achieved.
\end{exercise}

\subsection*{Derivation}\pathtag{standard}

\begin{exercise}
Derive the Paris law solution for crack-growth life with constant $\Delta K$ (which approximates short-crack behaviour or constant-amplitude loading on a slowly growing crack).
\end{exercise}

\begin{exercise}
Derive the Goodman relation for the effect of mean stress on fatigue strength.
\end{exercise}

\begin{exercise}
Derive the stress-intensity factor for a small surface crack in a uniformly stressed plate.
\end{exercise}

\begin{exercise}
Derive the relationship between stress concentration factor $K_t$ and fatigue notch factor $K_f$ for various notch sensitivities.
\end{exercise}

\subsection*{Estimation}\pathtag{standard}

\begin{exercise}
Estimate the fatigue life of a steel bracket experiencing $\sigma_a = 250\,\text{MPa}$, mean stress $100\,\text{MPa}$. State assumptions.
\end{exercise}

\begin{exercise}
Estimate the inspection interval for an aircraft fuselage that experiences $20{,}000$ pressurisation cycles per year, with a $2.5\,\text{mm}$ inspection-threshold flaw and a $25\,\text{mm}$ critical size.
\end{exercise}

\begin{exercise}
Estimate the fatigue-life reduction from surface roughness $R_a = 6.3\,\mu\text{m}$ compared to polished surface for a steel.
\end{exercise}

\begin{exercise}
Estimate the size of the plastic zone at a crack tip in a steel ($\sigma_y = 500\,\text{MPa}$) under $K = 60\,\text{MPa}\sqrt{\text{m}}$.
\end{exercise}

\subsection*{Simulation}\pathtag{standard}

\begin{exercise}
Simulate the rainflow counting of a synthetic load spectrum; apply Miner's rule and predict the life.
\end{exercise}

\begin{exercise}
Simulate crack growth under variable-amplitude loading with and without load-interaction effects (retardation).
\end{exercise}

\begin{exercise}
Simulate the probability of detection (POD) curve for an inspection technique; integrate over a crack-size distribution to compute the expected post-inspection flaw distribution.
\end{exercise}

\subsection*{Design}\pathtag{standard}

\begin{exercise}
Design a fatigue test program for a bracket: load spectrum, specimen geometry, instrumentation, sample size, target reliability.
\end{exercise}

\begin{exercise}
Design the damage-tolerance inspection program for an aircraft component: initial flaw assumption, growth analysis, intervals, NDI methods.
\end{exercise}

\begin{exercise}
Design a fail-safe structural design that ensures load redistribution after partial failure: redundant load paths, crack arrestors, leak-before-break.
\end{exercise}

\begin{exercise}
Design a fatigue-resistant detail for a welded steel structure: stress concentration mitigation, weld geometry, post-weld treatment.
\end{exercise}

\subsection*{Judgement}\pathtag{core}

\begin{exercise}
A team's component has cracked in service significantly before the predicted life. State three engineering hypotheses.
\end{exercise}

\begin{exercise}
A team is choosing between safe-life and damage-tolerance philosophy for a rotorcraft critical part. State the engineering trade-offs.
\end{exercise}

\begin{exercise}
A team's inspection program is missing cracks that propagate to failure. State the engineering response.
\end{exercise}

\begin{exercise}
A team is asked to extend the service life of a structural component beyond the original design life. State the engineering considerations.
\end{exercise}

\begin{exercise}
A team's load spectrum was based on assumed mission profiles; actual operation is more severe. State the engineering response.
\end{exercise}

\begin{exercise}
A team is choosing between LEFM and EPFM analysis for a structural integrity assessment. State the engineering criteria.
\end{exercise}

\begin{exercise}
A team's component has surface defects from manufacturing; the original fatigue analysis assumed defect-free material. State the engineering response.
\end{exercise}

\begin{exercise}
A team is asked to assess the safety of an aged welded structure with no documented fatigue analysis. State the engineering approach.
\end{exercise}

\begin{exercise}
A team is comparing two design alternatives for a fatigue-loaded component: one with traditional safety factor on stress, the other with damage-tolerance philosophy. State the engineering criteria for the choice, and discuss the additional considerations when the component will also operate in a corrosive environment not addressed in the original fatigue data.
\end{exercise}
